\documentclass{report}
\usepackage{amsmath,amssymb}
\setlength{\parindent}{0mm}
\setlength{\parskip}{1em}
\begin{document}
\begin{center}
\rule{6in}{1pt} \
{\large Lingfei Wu \\
{\bf Preconditioned Thick-Restart Lanczos Method with Subspace Optimization for Symmetric Eigenvalue Problems}}

120 Lafayette Blvd APT F \\ Williamsburg \\ VA \\ 23188
\\
{\tt lwu@email.wm.edu}\\
Fei  Xue\\
Andreas Stathopoulos\end{center}

The Thick-Restart Lanczos method (TRLAN) is an effective method for large
sparse symmetric eigenvalue problems. However, the effectiveness of
Lanczos method heavily depends on the distribution of eigenvalues and can
converge very slowly when the desired eigenvalues are not well separated
from the rest of the spectrum. In addition, restarting causes further
deterioration of convergence. Thick restarting can recover some of the
convergence loss but the restarting method of subspace optimization (also
known as +k or locally optimal restarting) is the one that can result in
near optimal methods. However, locally optimal restarting has only been
implemented with Generalized Davidson (GD) or LOPBCG methods which have
expensive iteration costs.

In this work, we propose a new near-optimal eigenmethod, TRLAN+K, that
combines the efficiency of the Lanczos algorithm and the power of the
subspace optimization techniques. Specifically, an inner-outer scheme is
employed. The inner loop applies the Lanczos recurrence while the outside
loop augments the Krylov subspace with a certain number of vectors from
the previous restart cycle. We then apply the Rayleigh-Ritz procedure in
the new subspace. Numerical experiments show that the proposed method can
achieve almost optimal convergence when considering the unrestarted
Lanczos as benchmark. Specifically, our method can converge up to ten
times faster than TRLAN when both methods are operated with relatively
small subspace under limited memory. The advantage of the proposed method
slowly diminishes as the size of subspace increases compared to TRLAN,
and then both methods converge similarly to the unrestarted Lanczos.

The generated Davidson (GD) method and its variants are a class of
effective eigenmethods that can also take advantage of a preconditioner.
Among them, GD+K utilizes the locally optimal conjugate gradient
restarting and has been shown to deliver near-optimal convergence under
severely limited memory. Without a preconditioner, it is well known that
the GD method is mathematically equivalent to the Lanczos method and with
thick restarting it is equivalent to TRLAN. GD+K, however, has slightly
better convergence than TRLAN+K since it expands the subspace by the
residuals from Rayleigh-Ritz optimization at every step. On the other
hand TRLAN+K has smaller costs for Raryleigh-Ritz and orthogonalization
and thus it can be more efficient than GD+K when the matrix is
sufficiently sparse. In addition, as the size of the subspace increases,
the differences between convergence rates become smaller and the low cost
of the Lanczos recurrence further benefits TRLAN+K.

In addition to the usual shift-invert mode,TRLAN+K can also be adapted to
use a preconditioner. We use an inverse-free type preconditioned Krylov
subspace scheme in the inner iteration. Although the inner iteration is
similar to the EIGIFP method, TRLAN+K differs in several ways and most
importantly in the techniques used in the outer iteration, specifically
thick restarting and the subspace optimization. These techniques are the
key to achieving a near-optimal eigenmethod.


\end{document}
